\subsection{Kebutuhan Fungsional}

Berdasarkan analisis permasalahan pada subbab~\ref{analisis-permasalahan} dan analisis solusi pada subbab~\ref{analisis-solusi}, dirancang kebutuhan fungsional pada Tabel~\ref{tab:kebutuhan-fungsional-pghd}.
Kebutuhan fungsional dikodekan dengan \verb|FXX|, di mana \verb|XX| adalah kode unik kebutuhan fungsional.

\begin{table}[htbp]
    \centering
    \caption{Kebutuhan fungsional integrasi PGHD pada DecMed}\label{tab:kebutuhan-fungsional-pghd}
    \begin{styledtable}{>{\raggedright\arraybackslash}p{0.08\textwidth} >{\raggedright\arraybackslash}p{0.10\textwidth} >{\raggedright\arraybackslash}X}
        \textbf{ID} & \textbf{Gap} & \textbf{Deskripsi Kebutuhan Fungsional} \\

        F01 & Gap 1 & Pasien dapat mengambil data PGHD secara otomatis dari perangkat (misalnya \textit{wearable} dan sensor) atau memasukkan data PGHD secara manual melalui antarmuka aplikasi klien. \\

        F02 & Gap 2 & Data PGHD yang dikumpulkan di perangkat pasien disimpan secara lokal sehingga tidak hilang meskipun perangkat dalam kondisi tanpa koneksi internet. \\

        F03 & Gap 2, Gap 3 & Sistem secara otomatis mengirimkan data PGHD yang sudah terkumpul ke server DecMed ketika koneksi internet tersedia, dengan mekanisme \textit{batch} yang mendukung pengiriman bertahap. \\

        F04 & Gap 1 & Sistem mengelompokkan dan menyajikan data PGHD berdasarkan sumber (perangkat \textit{wearable}, aplikasi seluler, dan laporan mandiri) sehingga dapat dibedakan oleh pasien maupun tenaga kesehatan. \\

        F05 & Gap 4 & Sistem menyimpan metadata PGHD yang mencakup waktu pengukuran, jenis sumber data, dan identitas pasien sebagai penghasil data untuk mendukung penelusuran asal-usul (\textit{data provenance}). \\

        F06 & Gap 5 & Pasien dapat memberikan izin akses baca PGHD kepada tenaga kesehatan untuk kunjungan tertentu melalui mekanisme kontrol akses yang terintegrasi dengan DecMed. \\

        F07 & Gap 5 & Tenaga kesehatan yang telah menerima izin akses dari pasien dapat melihat daftar entri PGHD yang terkait dengan pasien dan kunjungan yang sedang berlangsung. \\

        F08 & Gap 5 & Tenaga kesehatan dapat membuka dan melihat isi detail dari entri PGHD pasien yang dipilih, termasuk metadata asal-usul data (waktu, perangkat, dan jenis sumber). \\

        F09 & Gap 5 & Sistem menampilkan peringatan kepada tenaga kesehatan apabila data PGHD yang diakses terdeteksi telah rusak, dimodifikasi secara tidak sah, atau gagal melalui pemeriksaan integritas. \\

        F10 & Gap 5 & Pasien dapat mencabut izin akses terhadap data PGHD yang dimilikinya kapan saja, dan pencabutan tersebut langsung berlaku sehingga tenaga kesehatan tidak lagi dapat mengakses data terkait. \\
    \end{styledtable}
\end{table}